\subsection{Влажность}
%%%%%%%%%%%%%%%%%%%%%%%%%%%%%%%%%%%%%%%%%%%%%%%%%%%%
%%%%%%%%%%%%%%%%%%%%%%%%%%%%%%%%%%%%%%%%%%%%%%%%%%%%
\z{Какая относительная влажность была в сауне при температуре $+100\,\degree\text{C}$, если капли на стенах начали появляться при температуре $+70\,\degree\text{C}$.}
%
{$\varphi = \frac{P_{70}}{P_{100}} = 31.4\%$}

%%%%%%%%%%%%%%%%%%%%%%%%%%%%%%%%%%%%%%%%%%%%%%%%%%%%
\z{В цилиндре под поршнем находится воздух с относительной влажностью $\varphi = 70\%$. Объём цилиндра изотермически уменьшили в $10$ раз. Какая часть водяного пара сконденсировалась? Объёмом жидкости в конечном состоянии можно пренебречь.}
%
{$\varphi = 85.7\%$}

%%%%%%%%%%%%%%%%%%%%%%%%%%%%%%%%%%%%%%%%%%%%%%%%%%%%
\z{Определите массу $m$ росы, выпавшей при уменьшении объема воздуха в четыре раза, если начальный объем воздуха $V = 1\,\text{м}^3$, температура $t = 20\,\degree\text{C}$ и относительная влажность $\varphi = 50\%$. Температура все время постоянна, плотность насыщенных паров при данной температуре $\rho_{\text{нас}} = 1.7 \cdot 10^{-2}\,\text{кг/м}^3$.}
%
{$m = 4.32\,\text{г}$}

%%%%%%%%%%%%%%%%%%%%%%%%%%%%%%%%%%%%%%%%%%%%%%%%%%%%
\z{При изотермическом уменьшении объема воздуха в четыре раза образовалось $m = 10\,\text{г}$ воды. Определите относительную влажность воздуха в начальный момент, если начальный объем воздуха $V = 1\,\text{м}^3$, температура $t = 20\,\degree\text{C}$. Плотность насыщенных паров при данной температуре $\rho_{\text{нас}} = 1.7 \cdot 10^{-2}\,\text{кг/м}^3$.}
%
{$\varphi = 83.8\%$}

%%%%%%%%%%%%%%%%%%%%%%%%%%%%%%%%%%%%%%%%%%%%%%%%%%%%
\z{Под поршнем находится влажный воздух плотностью $\rho = 1.5\rho_{\text{нас}}$ при заданной температуре. При изотермическом сжатии в $4$ раза плотность стала равна $5\rho_{\text{нас}}$ при данной температуре. Чему была равна начальная относительная влажность?}
%
{$\varphi = 50\%$}

%%%%%%%%%%%%%%%%%%%%%%%%%%%%%%%%%%%%%%%%%%%%%%%%%%%%
\z{В цилиндре под поршнем находятся в равновесии воздух, водяной пар и вода. Отношение масс жидкости и пара $\alpha = \frac{1}{2}$. В медленном изотермическом процессе объём влажного воздуха увеличивается в $k = 3$ раза. Найти относительную влажность воздуха в цилиндре в конечном состоянии.}
%
{$\varphi = 50\%$}

%%%%%%%%%%%%%%%%%%%%%%%%%%%%%%%%%%%%%%%%%%%%%%%%%%%%
\z{В сосуде находится влажный воздух. При изотермическом сжатии его объём уменьшился в $5$ раз, а плотность увеличилась в $3$ раза. При дальнейшем изотермическом сжатии в $3$ раза плотность в итоге стала в $7$ раз больше первоначальной. Какую относительную влажность $\varphi$ имел воздух до начала сжатия?}
%
{$\varphi = 60\%$}


